\chapter{Related Work}

\section{Air Quality Research in Ukraine}

Air quality in Ukraine has been the subject of scientific investigation across multiple dimensions over the past decade. Pre-conflict studies focused on establishing baseline pollution levels and identifying key emission sources\cite{shelestov2020essential, svidzinska2020citizen}, while more recent work has examined the acute environmental impact of the ongoing military conflict\cite{MALYTSKA2026108959}. A persistent challenge across all of these efforts is the reliable collection of high-quality sensor data at sufficient spatial and temporal resolution.

Svidzinska\cite{svidzinska2020citizen} demonstrated that citizen-deployed sensor networks — comprising 133 sensors across 7 independent projects in Kyiv, collecting readings at 1–10 minute intervals — can serve as a meaningful complement or partial substitute for state-operated monitoring infrastructure, which remains sparse and technologically outdated\cite{pehchevski2020ukraine}. Shelestov et al.\cite{shelestov2020essential} proposed a conceptual infrastructure architecture for measuring PM2.5 and PM10 concentrations over Kyiv using satellite-derived data, ground-based observations, and atmospheric modelling. Zhurakovskyi et al.\cite{zhurakovskyi2025ml} proposed an IoT-based environmental monitoring system integrating machine learning for automated data collection and prediction of environmental state changes, with TinyML-based local inference to reduce dependence on cloud services.

Despite this body of work, the existing literature on Ukrainian air quality monitoring addresses individual components in isolation — sensor deployment, data collection, or predictive modelling — without integrating them into a pipeline architecture covering the full path from ingestion to alerting and visualization.

\section{Streaming Pipeline Technologies}
As documented by Carbone et al.\cite{carbone2020beyond}, the foundations of stream processing attracted the interest of database research communities as early as the 1990s. That decade witnessed the emergence of prototype systems such as TelegraphCQ and Stanford STREAM, which introduced foundational concepts including window aggregations, fault tolerance, and complex event processing\cite{carbone2020beyond}. Building on this research, the first commercial systems — including Esper and Oracle CQL — emerged in the 2000s. The 2010s marked a decisive shift toward widespread industrial adoption, with the introduction of Twitter Storm, Apache Spark Streaming, and Apache Flink, each targeting the growing demand for low-latency, large-scale data processing.

Stream data processing has since matured into one of the most widely adopted data processing paradigms. Modern use cases span IoT data aggregation, real-time fraud detection, and medical monitoring systems — all of which require minimal latency between data generation and actionable output. Kumar proposes a latency-based framework for architectural selection, noting that financial applications demand sub-50\,ms response times while behavioural analytics can tolerate latencies measured in minutes\cite{kumar2025evolution}. At the technology level, Apache Kafka and Apache Flink have emerged as the dominant combination for production streaming pipelines, with properly configured Kafka clusters sustaining throughput exceeding one million messages per second at sub-10\,ms latencies\cite{kumar2025evolution}.

The three frameworks most relevant to this work are compared below. \textbf{Apache Spark Streaming} extends the Spark batch engine with a micro-batch execution model: the incoming stream is divided into small time-bounded batches, each processed as a discrete RDD computation. This yields high throughput but introduces an inherent latency floor equal to the batch interval, typically 500\,ms to several seconds. \textbf{Apache Storm} was one of the first frameworks to offer true record-at-a-time processing, delivering low per-record latency. However, it provides no native support for stateful windowed aggregations or exactly-once processing guarantees, requiring these to be implemented manually at the application level. \textbf{Apache Flink} combines a true event-at-a-time execution model with native support for event-time semantics, watermarking, managed state, and exactly-once guarantees via distributed snapshots. Empirical benchmarks by Joy \cite{joy2024scalable} confirm that Flink achieves lower latency and higher throughput than both alternatives while consuming fewer resources. Table~\ref{tab:stream-frameworks} summarises the key properties of each framework.

\begin{table}[H]
\centering
\caption{Comparison of stream processing frameworks}
\label{tab:stream-frameworks}
\begin{tabular}{|l|l|l|l|}
\hline
\textbf{Property} & \textbf{Spark Streaming} & \textbf{Apache Storm} & \textbf{Apache Flink} \\
\hline
Execution model       & Micro-batch          & Record-at-a-time  & Event-at-a-time   \\\hline
Latency               & Seconds              & Low               & Very low          \\\hline
Managed state         & No                   & No                & Yes               \\\hline
Exactly-once          & Yes (micro-batch)    & No                & Yes               \\\hline
Native windowing      & Limited              & No                & Yes               \\\hline
Active development    & Yes                  & Minimal           & Yes               \\\hline
\end{tabular}
\end{table}

\section{Streaming Data Processing Pipelines in Practice}

The few publicly documented end-to-end pipeline implementations offer useful reference points for architectural decision-making. Kamal et al.\cite{kamal2024realtime} describe the deployment of a real-time visualization and analysis architecture at Cologne University Hospital's Medical Data Integration Center, leveraging the ELK stack for continuous ingestion and illustrating the applicability of streaming architecture in a regulated, data-intensive environment. Akanbi and Masinde\cite{akanbi2020distributed} propose a distributed stream processing middleware framework for real-time analysis of heterogeneous sensor data, demonstrating its application specifically in the environmental monitoring domain.

Kumar et al.\cite{kumar2024cep} propose a rule-based Complex Event Processing framework for air quality monitoring in smart cities, using Apache Kafka for stream ingestion and SPARQL queries over a knowledge graph for pattern-based alert detection. Sserunjogi et al.\cite{sserunjogi2025airqo} documented the AirQo data pipeline — a production cloud-native system for AI-driven air quality monitoring in Africa. The ingestion layer is orchestrated using Apache Airflow while BigQuery serves as the analytical backend. Kafka is employed as a message broker to decouple the ingestion layer from downstream microservices, enabling asynchronous, fault-tolerant data distribution. This work illustrates that real-world deployments often combine streaming and batch processing depending on the operational constraints of the environment.

\section{Summary}

Existing research on air quality in Ukraine has produced valuable contributions across sensor deployment, atmospheric modelling, and conflict-driven pollution analysis\cite{shelestov2020essential, svidzinska2020citizen, MALYTSKA2026108959}. However, these efforts address individual components in isolation and do not extend to integrated pipeline architectures capable of continuous ingestion, processing, and real-time notification.

The stream processing landscape has matured significantly, evolving from academic prototypes in the 1990s to production-grade frameworks such as Apache Flink and Apache Kafka that now underpin real-time systems across industries\cite{carbone2020beyond, kumar2025evolution}. Practical pipeline implementations have been documented across domains — from healthcare\cite{kamal2024realtime} and general environmental monitoring\cite{akanbi2020distributed} to air quality systems using CEP-based alerting\cite{kumar2024cep} and cloud-native ETL architectures\cite{sserunjogi2025airqo} — demonstrating that end-to-end streaming pipelines for sensor-driven domains are both feasible and publishable.

Taken together, these threads identify a clear gap: the availability of mature streaming technology and demonstrated domain need in Ukraine has not yet been bridged by a published pipeline architecture oriented toward sub-minute air quality notification. This work addresses that gap.